
\documentclass{article}
\usepackage{colortbl}
\usepackage{makecell}
\usepackage{multirow}
\usepackage{supertabular}

\begin{document}

\newcounter{utterance}

\centering \large Interaction Transcript for game `wordle', experiment `high\_frequency\_words\_no\_clue\_no\_critic', episode 6 with reflexion{-}Qwen3.5{-}9B.
\vspace{24pt}

{ \footnotesize  \setcounter{utterance}{1}
\setlength{\tabcolsep}{0pt}
\begin{supertabular}{c@{$\;$}|p{.15\linewidth}@{}p{.15\linewidth}p{.15\linewidth}p{.15\linewidth}p{.15\linewidth}p{.15\linewidth}}
    \# & \multicolumn{2}{c}{Player} && \multicolumn{2}{c}{Game Master} \\
    \hline

    \theutterance \stepcounter{utterance}  
    & & & \multicolumn{4}{p{0.6\linewidth}}{
        \cellcolor[rgb]{0.9,0.9,0.9}{
            \makecell[{{p{\linewidth}}}]{
                \texttt{\tiny{[P1$\langle$GM]}}
                \texttt{You have attempted this same task before and did not succeed. These are the plans you wrote after those attempts.} \\
\texttt{Trial \#1: New plan: I failed because I started generating a long, unformatted chain of thought and analysis instead of strictly adhering to the required "explanation:" and "guess:" format. Additionally, I incorrectly deduced the target word was "arbor" despite the feedback showing 'r' at index 1 was red (absent) while 'r' at index 4 was green (correct), which is a contradiction for the word "arbor" itself. My strategy of re{-}guessing "arise" was also redundant. Next time, I will immediately stop internal monologue, correctly parse the feedback (Target: \_ \_ \_ o r, contains 'a', 'o', 'r'; 'a' not at start; 'r' only at end; no 'b', 'i', 's', 'e'), and output a single, concise explanation followed by a valid five{-}letter word guess in the exact required format.} \\
\\ 
\texttt{Trial \#2: New plan: I failed because I misinterpreted the feedback from "aroma" (a<yellow> r<yellow> o<yellow> m<yellow> a<red>). The red 'a' at the end means the second 'a' is absent, but the yellow 'a' at the start means an 'a' exists elsewhere (likely index 3 or 4). I incorrectly assumed the word was "armor" based on a flawed deduction that 'a' wasn't at the start, ignoring that the first 'a' was yellow. I also failed to utilize the confirmed 'm' and the pattern of four yellows to deduce the specific permutation. Next time, upon seeing four yellows and one red for a 5{-}letter word, I will immediately calculate the only valid permutation where the red letter is removed and the others are shifted, rather than guessing a common word that might not fit the specific constraints. Specifically, for "aroma" feedback, I should have deduced the word is "armor" (wait, "armor" has 'a' at start and end? No, "armor" is a{-}r{-}m{-}o{-}r. Let's re{-}evaluate: Target has a, r, o, m. 'a' is not at start/end. 'r' is not at start. 'o' is not at start/end. 'm' is not at start. Actually, the target is likely "armor" but my logic was messy. Let's look at the actual target: The feedback for aroma was a(y) r(y) o(y) m(y) a(r). This means the set is \{a, r, o, m\} plus one more letter. The 'a' at the end is wrong, so the 'a' is elsewhere. The 'r', 'o', 'm' are in wrong spots. The correct word is likely "armor" if the target was a{-}r{-}m{-}o{-}r? No, "armor" has 'a' at start. If target is "armor", guess "aroma": a(y) r(y) o(y) m(y) a(r). This fits if target is a{-}r{-}m{-}o{-}r? No. Let's try "armor" as target. Guess "aroma": a(y) r(y) o(y) m(y) a(r). This implies 'a' is in the word but not at pos 0 or 4. 'r' is in word but not at 1. 'o' is in word but not at 2. 'm' is in word but not at 3. This is consistent with "armor" (a{-}r{-}m{-}o{-}r)? No, 'r' is at 1 in "armor". Wait, the feedback said r<yellow>. So 'r' is not at 1. So "armor" is wrong. The word must be "roam" + something? Or "} \\
\\ 
\\ 
\texttt{You are a language wizard who likes to guess words by using the given rules.} \\
\\ 
\texttt{Welcome to Wordle! You have six attempts to guess the target word, a valid English word of five lowercase letters (a{-}z). Please use the tags "explanation:" and "guess:" to provide a concise explanation for each guess.} \\
\\ 
\texttt{For instance, if your guess is "apple", your response should be} \\
\texttt{explanation: this is a common five{-}letter English word, and I am starting my guess with this word.} \\
\texttt{guess: apple} \\
\\ 
\texttt{After each guess, your answer will be validated, and you will receive feedback indicating which letters are correct (green), which letters are correct but in the wrong position (yellow), and which letters are incorrect (red). This feedback can be useful in determining which letters to include or exclude in your next guess.} \\
\\ 
\texttt{For example, the feedback for "apple" might be:} \\
\texttt{guess\_feedback: a<yellow> p<yellow> p<green> l<yellow> e<red>} \\
\\ 
\texttt{The explanation should contain details about how the guess\_feedback is used to arrive at a new guess.} \\
\\ 
\texttt{Do not generate any other text.} \\
\\ 
\texttt{Let's begin with your first guess.} \\
            }
        }
    }
     \\ \\

    \theutterance \stepcounter{utterance}  
    & \multicolumn{4}{p{0.6\linewidth}}{
        \cellcolor[rgb]{0.561,0.737,0.561}{
            \makecell[{{p{\linewidth}}}]{
                \texttt{\tiny{[P1$\rangle$GM]}}
                \texttt{explanation: Starting with a common five{-}letter word that contains frequent letters to gather initial feedback.} \\
\texttt{guess: apple} \\
            }
        }
    }
    & & \\ \\

    \theutterance \stepcounter{utterance}  
    & & & \multicolumn{2}{p{0.3\linewidth}}{
        \cellcolor[rgb]{0.9,0.9,0.9}{
            \makecell[{{p{\linewidth}}}]{
                \texttt{\tiny{[GM$|$GM]}}
                \texttt{attempts = 1} \\
\texttt{target\_word = major} \\
\texttt{guess = apple} \\
\texttt{guess\_feedback = a<yellow> p<red> p<red> l<red> e<red>} \\
            }
        }
    }
    & & \\ \\

    \theutterance \stepcounter{utterance}  
    & & & \multicolumn{4}{p{0.6\linewidth}}{
        \cellcolor[rgb]{0.9,0.9,0.9}{
            \makecell[{{p{\linewidth}}}]{
                \texttt{\tiny{[P1$\langle$GM]}}
                \texttt{guess\_feedback: a<yellow> p<red> p<red> l<red> e<red>} \\
\\ 
\texttt{Provide your response only in this format.} \\
\texttt{explanation: details} \\
\texttt{guess: word} \\
            }
        }
    }
     \\ \\

    \theutterance \stepcounter{utterance}  
    & \multicolumn{4}{p{0.6\linewidth}}{
        \cellcolor[rgb]{0.561,0.737,0.561}{
            \makecell[{{p{\linewidth}}}]{
                \texttt{\tiny{[P1$\rangle$GM]}}
                \texttt{explanation: The feedback 'a<yellow>' indicates that 'a' is in the word but not at the first position, while 'p', 'l', and 'e' are absent. I will try a new word starting with 's' and including 'a' in a different position to narrow down the possibilities.} \\
\texttt{guess: starry} \\
            }
        }
    }
    & & \\ \\

    \theutterance \stepcounter{utterance}  
    & & & \multicolumn{2}{p{0.3\linewidth}}{
        \cellcolor[rgb]{0.9,0.9,0.9}{
            \makecell[{{p{\linewidth}}}]{
                \texttt{\tiny{[GM$|$GM]}}
                \texttt{Error: The length of the guessed word is not 5.} \\
            }
        }
    }
    & & \\ \\

    \theutterance \stepcounter{utterance}  
    & & & \multicolumn{2}{p{0.3\linewidth}}{
        \cellcolor[rgb]{0.9,0.9,0.9}{
            \makecell[{{p{\linewidth}}}]{
                \texttt{\tiny{[GM$|$GM]}}
                \texttt{game\_result = ABORT} \\
            }
        }
    }
    & & \\ \\

    \theutterance \stepcounter{utterance}  
    & & \multicolumn{4}{p{0.6\linewidth}}{
        \cellcolor[rgb]{0.95,0.95,0.95}{
            \makecell[{{p{\linewidth}}}]{
                \texttt{\tiny{[P1$\leftrightarrow$P]}}
                \texttt{Post{-}episode Reflexion (not sent to GM):} \\
\texttt{New plan: I failed because I generated a guess ("starry") that is six letters long, violating the strict rule of using only five lowercase letters. Additionally, my explanation incorrectly claimed 'p' and 'l' were absent when the feedback `p<red> l<red>` actually confirmed they are not in the word, but I failed to account for the fact that "starry" contains two 'r's and two 'y's, which might not be optimal given the constraints. Next time, I will immediately verify that my guess is exactly five characters long before outputting it. Based on the feedback `a<yellow> p<red> p<red> l<red> e<red>`, I know the word contains 'a' (not at index 0) and excludes 'p', 'l', 'e'. I will construct a valid five{-}letter word like "actor" or "area" (checking constraints) and ensure the output strictly follows the `explanation:` and `guess:` format with no extra text. Specifically, I will choose a word that places 'a' in a non{-}first position and uses common consonants to maximize information gain, ensuring the word length is exactly 5.} \\
            }
        }
    }
    & \\ \\

\end{supertabular}
}

\end{document}
